\centerline{\bf \Huge Счёт в синусах}

\begin{flushright}
        \scriptsize{Эпизод 8. Аска прибывает в Японию}
\end{flushright}

\textbf{Памятка}

$\sin (-\alpha)=-\sin \alpha, \cos (-\alpha)=\cos \alpha,\quad \sin ^2 \alpha+\cos ^2 \alpha=1$;

$\sin \left(\frac{\pi}{2}-\alpha\right)=\cos \alpha,\  \cos \left(\frac{\pi}{2}-\alpha\right)=\sin \alpha,\  \sin (\pi-\alpha)=\sin \alpha, \ \cos (\pi-\alpha)=-\cos \alpha ;$

$\sin (\alpha \pm \beta)=\sin \alpha \cos \beta+\sin \beta \cos \alpha,\quad \cos (\alpha \pm \beta)=\cos \alpha \cos \beta \mp \sin \alpha \sin \beta ;$ 

$\sin 2 \alpha=2 \sin \alpha \cos \alpha, \quad \cos 2 \alpha=2 \cos ^2 \alpha-1=1-2 \sin ^2 \alpha$.

\textbf{Теорема Чевы в синусной форме.} Дан треугольник $A B C$. Прямые $A A_1, B B_1, C C_1$ пересекаются в одной точке тогда и только тогда, когда

$$
\frac{\sin \angle\left(A B, A A_1\right)}{\sin \angle\left(A A_1, A C\right)}=\frac{\sin \angle\left(C A, C C_1\right)}{\sin \angle\left(C C_1, C B\right)} \cdot \frac{\sin \angle\left(B C, B B_1\right)}{\sin \angle\left(B B_1, B A\right)}=1 .
$$


\problem{1}
В треугольнике $A B C$ провели (а) высоты (б) биссектрисы. Выразите длины всех получившихся отрезков через $\angle A=\alpha, \angle B=\beta, \angle C=\gamma$ и радиус описанной окружности $R$.

\problem{2}
В остроугольном треугольнике $A B C$ проведены высоты $A A_1$ и $B B_1$. Точка $O$ - центр описанной окружности. Докажите, что расстояние от точки $A_1$ до прямой $B O$ равно расстоянию от точки $B_1$ до прямой $AO$.

\problem{3}
В треугольнике с тупым углом $A$ проведены высоты $B B_1$ и $C C_1$. Докажите, что отрезок, соединяющий проекции точки $B_1$ на прямые $B A$ и $B C$ равен отрезку, соединяющему проекции точки $C_1$ на прямые $CA$ и $CB$.

\problem{4}
Внутри угла с вершиной $O$ отмечена точка $P$. Рассматриваются всевозможные пары точек $X$ и $Y$ на сторонах угла такие, что $\angle O P X=\angle O P Y$. Докажите, что все прямые $X Y$ проходят через одну точку.

\problem{5}
Точка $X$ лежит внутри правильного треугольника $A B C$. Точки $A_1, B_1, C_1$ симметричны точке $X$ относительно сторон $B C, A C, A B$ соответственно. Докажите, что прямые $A A_1, B B_1, C C_1$ пересекаются в одной точке.

\newpage

\hspace{3mm}

\problem{6}
Две окружности радиусов $r$ и $R$ касаются прямой в точках $A$ и $B$. Пусть $C$ - одна из точек пересечения этих окружностей. Докажите, что радиус описанной окружности треугольника $A B C$ не зависит от положения окружностей.

\problem{7}
В остроугольном треугольнике $A B C$ проведены высоты $A A_1, B B_1, C C_1$. Докажите, что треугольник с вершинами в ортоцентрах треугольников $A B_1 C_1, B C_1 A_1, C A_1 B_1$ равен треугольнику $A_1 B_1 C_1$.

\problem{8}
Две окружности касаются друг друга внутренним образом в точке $N$. Касательная к внутренней окружности, проведённая в точке $K$, пересекает внешнюю окружность в точках $A$ и $B$. Пусть $M$ - середина дуги $A B$, не содержащей точку $N$. Докажите, что радиус описанной окружности треугольника $B M K$ не зависит от выбора точки $K$ на внутренней окружности.

\problem{9}
На сторонах $A B, B C, C A$ треугольника $A B C$ выбраны точки $C_1, A_1, B_1$ соответственно. Оказалось, что центры вписанных окружностей треугольников $A B C$ и $A_1 B_1 C_1$ совпадают, а радиус вписанной окружности в треугольник $A_1 B_1 C_1$ вдвое меньше, чем радиус окружности треугольника $A B C$. Докажите, что треугольник $A B C$ правильный.